\chapter{Soluciones de los casos de teoría avanzada}
\label{ch:matmo-teoria-avanzada-sol}
\noindent Las respuestas modelo señalan el núcleo mínimo de un razonamiento correcto; se valoran respuestas alternativas técnicamente justificadas dentro de la normativa aplicable.

\section*{AT01. Fractura y geometría}
$a=4,mathrm{mm}$ en esa convención. La geometría modifica el campo de intensidad y el factor adimensional. La tenacidad debe corresponder al material, temperatura, espesor y condición de carga; también debe examinarse la hipótesis elástica y el tamaño de la zona plástica.

\section*{AT02. Fatiga y tensión media}
Los valores de amplitud y media difieren: $sigma_a=(sigma_{max}-sigma_{min})/2$, $sigma_m=(sigma_{max}+sigma_{min})/2$. Para $R=-1$ la media es cero y la amplitud coincide con el máximo; para $R=0$ son ambas la mitad. Curvas, correcciones de media, entorno y dispersión deben ser compatibles.

\section*{AT03. Paris y predicción de inspecciones}
La integración estima ciclos de crecimiento en un intervalo bajo un espectro/modelo definidos. Faltan iniciación y regiones fuera de Paris, dispersión y factores de seguridad, detección de defectos, criterios de aceptación y la fuente regulatoria o de diseño. Un número calculado no es automáticamente un intervalo de mantenimiento.

\section*{AT04. Fluencia y álabes monocristalinos}
Eliminar límites internos puede reducir contribuciones de deslizamiento, difusión o cavitación intergranular, pero no la actividad de dislocaciones, difusión intragranular, fatiga o oxidación. La selección requiere temperatura, tensión y microestructura concretas; el monocristal no cambia automáticamente el punto de fusión.

\section*{AT05. Material compuesto orientado}
El refuerzo alineado domina la rigidez longitudinal; en dirección transversal participa de otro modo la matriz y la interfase. La regla longitudinal presupone compatibilidad de deformaciones; la inversa transversal es idealizada. Orientación de capas y secuencia de apilado afectan la respuesta global, incluida la matriz ABD.

\section*{AT06. Selección END}
Los penetrantes requieren abertura a superficie; las partículas magnéticas, ferromagnetismo. Para ultrasonidos y radiografía importan material, espesor, geometría, orientación, acceso, calibración, sensibilidad/POD y procedimiento aplicable. Una inspección negativa es evidencia condicionada por su cobertura.

\section*{AT07. Diseño, aeronave y mantenimiento}
El certificado de tipo se refiere al diseño aprobado, el certificado de aeronavegabilidad a la aeronave individual y su condición de validez, y el CRS a la liberación de trabajos realizados dentro del marco aplicable. Ninguno sustituye automáticamente los demás; deben comprobarse configuración, cumplimiento y registros.

\section*{AT08. Parte 145, 66 y 147}
La formación, licencia, habilitación de tipo/limitaciones cuando correspondan y autorización certificadora en una organización son escalones distintos. La Parte 145 define el alcance y sistema de la organización; ni formación ni licencia solas amplían privilegios organizativos. Depende del trabajo, de la norma aplicable y de autorización demostrable.

\section*{AT09. AD y boletín de servicio}
Se consulta la AD vigente, su aplicabilidad, revisiones y métodos que acepte, plazos y posibles medios alternativos de cumplimiento aprobados. El número de revisión mayor no implica equivalencia jurídica ni técnica, y debe conservarse evidencia del cumplimiento requerido.

\section*{AT10. ARC con hallazgos y plazos}
EASA aclara que no se emite ARC o recomendación con hallazgos abiertos; deben corregirse y verificarse. La revisión y la extensión son actos distintos con supuestos y plazos propios; las previsiones de 90 días para determinadas revisiones físicas y 30 días para anticipar ciertas extensiones no dan una tolerancia general de mantenimiento.

\section*{AT11. Tarea con dos límites y LLP}
Los vencimientos se evalúan con el registro de última ejecución y ambas bases. El margen de FH no anula el límite de calendario. Si la LLP carece de ciclos totales trazables, se mantiene estado indeterminado y se solicitan registros; un campo vacío no equivale a cero.

\section*{AT12. EGT, vibración y diagnóstico}
No de forma automática. La mejora de EGT puede apoyar una hipótesis, pero no excluye daño independiente. Deben normalizarse régimen, ambiente y sensores; relacionar la vibración con rotor de referencia, analizar partículas, inspeccionar según datos aplicables y documentar hallazgos y cierre por personal autorizado.

\noindent\textbf{Fuentes de control:} manual MATMO, capítulos 2--6, 9, 11 y 14--22; MIT OCW sobre mecánica de materiales; FAA AMT General; Reglamento UE 1321/2014; EASA Continuing Airworthiness y FAQ oficial de revisión ARC. Toda respuesta sobre aeronaves reales exige comprobar norma, versión, ámbito y documentación de mantenimiento.
